% Source: https://artofproblemsolving.com/community/c3h2143808_chess_problem
% Stable identity: aops-topic-2143808-q01
Number the rows and columns $1$ to $8$. A bishop on $(r,c)$ moves along the diagonals on which $r-c$ is constant or $r+c$ is constant. If two bishops share a diagonal, then the two closest bishops on that diagonal attack each other, so ``no two attack each other'' means exactly that no two bishops share a diagonal of either kind.

\textbf{Splitting by colour.} Bishops on squares of different colours (different parity of $r+c$) never share a diagonal. So if $a_k$ is the number of ways to place $k$ non-attacking bishops on the squares with $r+c$ even, and $b_k$ the same for $r+c$ odd, the answer is $\sum_k a_kb_{8-k}$. The reflection $c\mapsto9-c$ maps the board to itself, preserves diagonals, and swaps the two colours (because $9$ is odd), so $b_k=a_k$.

\textbf{One colour as a rook problem.} For the squares with $r+c$ even, put $u=r-c$ and $v=r+c$. The values of $u$ are $0,\pm2,\pm4,\pm6$, and the diagonal with a given $u$ contains the squares with
\begin{itemize}
\item $u=0$: $v\in\{2,4,\dots,16\}$ ($8$ squares);
\item $u=\pm2$: $v\in\{4,6,\dots,14\}$ ($6$ squares each);
\item $u=\pm4$: $v\in\{6,8,10,12\}$ ($4$ squares each);
\item $u=\pm6$: $v\in\{8,10\}$ ($2$ squares each).
\end{itemize}
Placing non-attacking bishops means choosing squares with all $u$ different and all $v$ different, that is, placing non-attacking rooks on a board with columns $u$ and rows $v$, where the columns have heights $2,2,4,4,6,6,8$ and each shorter column's set of rows is contained in each longer one's.

\textbf{Counting.} Add the columns in order of increasing height. If $j$ rooks are already in shorter columns, they occupy $j$ different rows, all of which lie in the new column of height $h$, so a rook can be added in the new column in $h-j$ ways. Hence, if $R(k)$ is the number of placements of $k$ rooks so far, adding a column of height $h$ changes it to
\[R'(k)=R(k)+(h-k+1)\,R(k-1).\]
Starting from $R(0)=1$:
\[\begin{array}{c|rrrrrrrr}
\text{heights used}&k=0&1&2&3&4&5&6&7\\ \hline
2&1&2&&&&&&\\
2,2&1&4&2&&&&&\\
2,2,4&1&8&14&4&&&&\\
2,2,4,4&1&12&38&32&4&&&\\
2,2,4,4,6&1&18&98&184&100&8&&\\
2,2,4,4,6,6&1&24&188&576&652&208&8&\\
2,2,4,4,6,6,8&1&32&356&1704&3532&2816&632&16
\end{array}\]
So $(a_0,a_1,\dots,a_7)=(1,32,356,1704,3532,2816,632,16)$, and $a_8=0$ (at most $7$ bishops fit on one colour, since there are only $7$ values of $u$).

\textbf{Combining.} The number of placements of $8$ bishops is
\[\sum_{k=1}^{7}a_ka_{8-k}=2\left(a_1a_7+a_2a_6+a_3a_5\right)+a_4^2=2\left(512+224992+4798464\right)+12475024=\boxed{22522960}.\]
